% !TeX root = ../main.tex
\chapter{实验与结果展示示例}
\label{chap:experiment-example}

本章展示普通表格和复杂结果表的写法。正式论文中，实验章节应包括实验目标、设置、数据、指标、结果、分析和小结。没有真实实验日志时，不要填写虚构数值。

\section{实验设置}
\label{sec:experiment-setup}

表题格式由模板自动设置为“表章号-表号  表题”。普通表格默认使用五号字、1.5 倍行距和三线表风格。

\begin{table}[htbp]
  \centering
  \caption{实验设置字段示例}
  \label{tab:setup-example}
  \begin{tabular}{C{0.25\textwidth}P{0.55\textwidth}}
    \toprule
    项目 & 填写说明 \\
    \midrule
    数据集 & 写明数据来源、版本、样本数和划分方式。 \\
    模型 & 写明模型名称、版本、参数规模和来源。 \\
    指标 & 写明指标定义、统计单位和阈值选择方式。 \\
    \bottomrule
  \end{tabular}
\end{table}

\section{结果表格}
\label{sec:result-table}

复杂结果表可以调用 \verb|\nbucompacttable|，该命令只作用于下一张表格，适合列较多的实验结果。示例见表~\ref{tab:compact-result-example}。

% !TeX root = ../main.tex
% 复杂实验结果表例外：列较多时可在 table 内调用 \nbucompacttable，使表内字号和行距局部收紧。
\begin{table}[htbp]
  \centering
  \caption{复杂结果表占位示例}
  \label{tab:compact-result-example}
  \nbucompacttable
  \begin{tabular}{C{0.20\textwidth}C{0.16\textwidth}C{0.16\textwidth}C{0.16\textwidth}C{0.16\textwidth}}
    \toprule
    方法 & 指标一 & 指标二 & 指标三 & 备注 \\
    \midrule
    Baseline & -- & -- & -- & 替换为真实结果 \\
    Ours & -- & -- & -- & 替换为真实结果 \\
    \bottomrule
  \end{tabular}
\end{table}

